\subsection{The triangular two-object exclusion through fourteen arrows}

\begin{lemma}[At most two elements outside the minimal ideal]
\label{C:two-outside-act-obstruction}
Let $M$ be a finite monoid with minimal ideal $K$, and suppose
$|M\setminus K|\leq2$.  For every finite-dimensional right
$R=\mathbb QM$-module $L$, eventual vanishing of
$\operatorname{Tor}_n^R(L,\mathbb Q)$ implies its vanishing for
every $n\geq3$.  In particular no finite right-act cone over $M$
has type $\mathbb{CP}^2$.
\end{lemma}
\begin{proof}
If $K=M$, the monoid is a finite group and its rational algebra
is semisimple.  Otherwise put $A=R/\mathbb QK$.
The Rees factorization
$\mathbb QK\cong Re\otimes_{\mathbb QH}eR$, with $eR$ free on
the left over the maximal group algebra, shows that $\mathbb QK$
is projective on the left over $R$.  Hence
$\operatorname{pd}_R A\leq1$.
Averaging $H$ gives a projective left module $P$ with its row-set
basis and an exact sequence
\[
0\longrightarrow V\longrightarrow P\longrightarrow\mathbb Q
\longrightarrow0.
\]
Every element of $K$ acts constantly on the row set, so it
annihilates the augmentation module $V$.  Thus $V$ is an $A$-module.

A rational algebra of dimension at most two is semisimple or is
$\mathbb Q[\varepsilon]/(\varepsilon^2)$.  In the semisimple case
$V$ is projective over $A$, so its projective dimension over $R$
is at most one.  This gives the asserted vanishing directly.
In the remaining case, Jordan normal form for the square-zero
operator gives
$V\cong A^a\oplus\mathbb Q^b$ as an $A$-module.
For $n\geq3$, dimension shifting therefore identifies
\[
\operatorname{Tor}_n^R(L,\mathbb Q)
 \cong\operatorname{Tor}_{n-1}^R(L,\mathbb Q_A)^b,
\]
where $\mathbb Q_A$ is the simple dual-number module restricted
to $R$.
If $b=0$ the conclusion follows.  If $b>0$, eventual vanishing
of these groups and duality give eventual vanishing of
$\operatorname{Ext}_R^j(\mathbb Q_A,DL)$.
The derived adjunction in Lemma~\ref{C:derived-Frobenius-Tor}
identifies these groups with
\[
\operatorname{Ext}_A^j\bigl(\mathbb Q,
                     \operatorname{RHom}_R(A,DL)\bigr).
\]
The second argument has cohomology only in degrees zero and one.
The minimal-injective-complex argument proved in that lemma
therefore makes these groups zero for every $j\geq2$.
Dimension shifting proves the claim.  Finally apply
Lemma~\ref{C:act-cone-Tor}; a projective-plane cone would have
nonzero $\operatorname{Tor}_3$ and an eventually zero tail.
\end{proof}

\begin{lemma}[A two-arrow biset opposite the sphere monoid]
\label{C:two-arrow-sphere-opposite}
Let $P=\{1\}\sqcup(\{0,1\}\times\{0,1\})$ be the
five-element sphere monoid.  If $H^2(BM;\mathbb Q)=0$, no
triangular biset category with vertex monoids $M,P$ and two
cross arrows has the rational cohomology ring of
$\mathbb{CP}^2$.  This requires no rational acyclicity of $BM$.
\end{lemma}
\begin{proof}
Pass to the opposite category if needed and write the biset as
a left $P$-set and right $M$-set.  Every element of the minimal
band of $P$ is idempotent.  Its image in the transformations of
a two-point set is therefore either the identity or a constant.
The rectangular-band equations imply, up to permuting the two
states, exactly three types: every band element acts identically;
every band element is the same constant; or the two rows act by
the two distinct constants.  Mixing an identity and a constant
would give a nontrivial semilattice image, which a rectangular
band cannot have.

In the first type the edge bar is $E_S\simeq E_R\times BP$,
and its map to $BP$ is projection.  Pullback on $H^2(BP;\mathbb Q)$
is injective.  Since $H^2(BM;\mathbb Q)=0$, the vertex restriction
of every degree-two class on the collage is zero.  Such classes
come from the relative cohomology of the two vertex spaces;
their products vanish because the corresponding relative quotient
is a suspension of the edge bar with a disjoint basepoint.

In the common-constant type, commutation forces the constant state
to be fixed by $M$ as well.  This gives a section
$BM\times BP\to E_S$ of the edge projection.  The map from
the direct sum of the two vertex $H^2$ groups into the edge group
is injective, so the same relative argument applies.

In the last type, a transformation commuting with both constants
fixes both states, so the $M$-action is trivial.  The left $P$-set
is its two-element minimal left ideal.  Its rational permutation
module is projective, and its action bar is rationally acyclic.
Thus $E_S\to BM$ is a rational homology equivalence.  The
homotopy-pushout description makes $BP\to BC$ a rational homology
equivalence, and in particular $H^4(BC;\mathbb Q)=0$.
All three cases exclude the stated ring.
\end{proof}

\begin{proposition}[The ten-element three-row act calculation]
\label{C:ten-three-row-three-point-acts}
No act-cone category with a ten-element vertex monoid and three
cross arrows has type $\mathbb{CP}^2$.
\end{proposition}
\begin{proof}
A vertex with zero rational $H^2$ is excluded by the relative-cup
argument.  A vertex with at most two minimal rows is excluded by
Theorem~\ref{C:two-row-act-cups}.  Of the remaining minimal-ideal
shapes at order ten, $4\times2$ and $3\times3$ bands have at
most two elements outside the ideal, and
Lemma~\ref{C:two-outside-act-obstruction} applies.  A nontrivial
minimal group with at least two rows and columns must have order
two and shape $2\times2$, already covered by the two-row theorem.
Thus only a $3\times2$ rectangular-band ideal remains.

The complete translation-table enumeration
\path{code/b_ten_non_square_targets.cpp} places the identity at
zero, the six band elements at $1,\ldots,6$, and the three other
elements at $7,8,9$.  There are $3^3 2^2=108$ possible linked
row/column transformation pairs.  Every multiset of three pairs
is tried, giving $\binom{110}{3}=215820$ profiles.  For every
ordered pair of outside elements, the possible product labels
are exactly those with the composed transformation pair.
All these labels are branched over, with associativity propagation;
every completed table is checked on every triple.  Isomorphism
uses every row, column, and outside permutation.  Since the
minimal ideal is characteristic, this is a complete enumeration.
It gives $58187$ associative sorted tables and $4260$ ordinary
isomorphism classes.  The opposite $2\times3$ enumeration gives
the same class count and is retained as a separate control.

Exact balanced ranks discard $1676$ of the $3\times2$ tables.
Another $1276$ have independently checked integral projective
resolutions of length at most two, so every act has zero fourth
rational homology.  For each of the remaining $1308$ tables,
all right actions on three states are enumerated by incremental
multiplication propagation, with exact stabilizers removing only
state conjugates.  There are $22741$ action classes.  The
disjoint rejection counts are
\[
\begin{array}{l|r}
\text{disconnected act, hence nonzero cone }H_1&5457\\
\text{wrong rational homology in degrees one through three}&14251\\
\text{wrong mod-two second Betti number}&88\\
\text{explicit mod-two primitive for the degree-two square}&2936\\
\text{nonzero square but nonzero mod-two }H_3&9.
\end{array}
\]
All nine last cases have $(b_3,b_4)=(1,1)$ over $\mathbb F_2$.
The rational calculations tensor an exact projective prefix through
$P_3$ with the act augmentation module, which computes
$\operatorname{Tor}_0,\operatorname{Tor}_1,\operatorname{Tor}_2$;
no omitted tail is used.  Every prefix is checked for literal
equivariance, zero composites, integral saturation, and exact rank
sums.  No resolution search failed on this list.
The normalized-bar primitives are checked on every degree-four
simplex.  Tables, actions, prefixes, primitives, and residual ranks
are in \path{code/outputs/b_ten_non_square_acts.json} and its
referenced compressed certificate chunks.
\end{proof}

\begin{proposition}[The triangular two-object exclusion through fourteen]
\label{C:fourteen-triangular}
No triangular finite category with two objects and at most fourteen
arrows has classifying space homotopy equivalent to
$\mathbb{CP}^2$.
\end{proposition}
\begin{proof}
The thirteen-arrow theorem handles smaller totals.  A singleton
cross-biset is a join and has zero positive-degree products.
For two cross arrows, Lemma~\ref{C:two-object-cup-pruning} leaves
the vertex-order pairs $(7,5)$ and $(6,6)$.  In the first pair,
Lemma~\ref{C:two-arrow-sphere-opposite} handles a seven-element
vertex with zero rational $H^2$; thus a possible $S^3$ vertex is
not silently called acyclic.  In the second pair, every six-element
vertex with zero rational $H^2$ is rationally acyclic by the
complete small-monoid classification.  It remains to enumerate
the displayed two-arrow profiles with both vertices having
nonzero rational $H^2$.

For at least three cross arrows, some vertex must have nonzero
rational $H^2$.  Choose it as the first vertex, passing to an
opposite category when appropriate.  If the other vertex has
order one or two, its rational coinvariants are exact.
Theorem~\ref{C:nine-act-cones}, the actual orbit-act comparison,
and Proposition~\ref{C:ten-three-row-three-point-acts} handle
all these cases.  The remaining profiles have both vertex orders
at least three and appear in the table below.  The additional
$(5,6,3)$ row covers the case where only the smaller vertex has
nonzero $H^2$.

The first vertex ranges over the complete lists of one, five,
forty-three, or $516$ possible $H^2$-positive tables at orders
five through eight, retaining the harmless zero-rank entries
in the order-eight list.  The other vertex ranges over all
ordinary classes of its order, except in the two-arrow rows,
where the preceding reductions permit both $H^2$-positive lists.
Every other-vertex monoid with projective trivial right rational
module is discarded using the all-act theorem through order nine;
each such discard is certified by an actual vector $v$ with
$vm=v$ for all $m$ and augmentation one.

For each remaining right act $S$, the program first enumerates
its complete endomorphism monoid by assigning the image of a
state and propagating equivariance on every monoid element.
It then enumerates every homomorphism from the opposite of the
other vertex monoid into this endomorphism monoid.  All power
and multiplication relations are checked.  No state permutation
is removed at this second stage, so every relative position of
the left action is retained.  This proves completeness of the
commuting-biset construction.  The disjoint checks give
\[
\begin{array}{c|r|r|r|r|r}
(m,n,s)&\text{bisets}&b_1\ne0&b_2\ne1&x^2=0&x^2\ne0\\\hline
(8,3,3)&43270&6636&24619&11985&30\\
(7,3,4)&30986&6400&13217&11343&26\\
(6,3,5)&27041&7181&13303&6557&0\\
(5,3,6)&55095&16503&19929&18663&0\\
(7,4,3)&25357&3782&13533&8014&28\\
(6,4,4)&27222&5743&13724&7735&20\\
(5,4,5)&67385&15627&36291&15467&0\\
(6,5,3)&23888&3773&12342&7743&30\\
(5,5,4)&55739&12078&26440&17212&9\\
(5,6,3)&53201&8623&26472&18076&30\\
(7,5,2)&430&43&129&257&1\\
(6,6,2)&254&25&79&135&15.
\end{array}
\]
All Betti numbers in this table are over $\mathbb F_2$.
Every one of the $189$ nonzero-square cases has nonzero third
mod-two homology.  Each zero square has an explicit normalized
degree-three primitive, checked on all degree-four simplices.
The complete run is \path{code/b_fourteen_small_bisets.py}, with
summary \path{code/outputs/b_fourteen_small_bisets.json} and
twelve referenced compressed files.  These contain every vertex
table, every projectivity splitting used to skip a vertex, every
primitive, and every residual homology calculation.  No numerical
tail vanishing is inferred.  All possible profiles are therefore
excluded.
\end{proof}

\begin{remark}\label{C:fourteen-triangular-verification-scope}
The above table and the ten-element calculation are complete exact
runs of their stated enumeration and cochain algorithms.  The additional
independent biset gates, which were still running at the ninth checkpoint,
are now complete.  The fresh pointwise-domain constraint solver
\path{code/b_verify_fourteen_biset_coverage.cpp} uses neither the
producer's monoid generators nor its endomorphism-monoid enumeration.
It independently enumerates $185847$ labelled right actions, passes
through $1176994$ constraint-search nodes, and reconstructs all
$409868$ commuting bisets.  Every per-vertex first- and second-homology
count agrees, and there is an exact bijection between its $123376$
degree-two survivors and the saved certificate records.  Its complete
output is \path{code/outputs/b_fourteen_biset_coverage_verified.jsonl}.

The separate verifier
\path{code/b_verify_fourteen_biset_primitives.cpp} reconstructs every
category from its two monoids and two actions.  It checks $73075962$
typed associativity equations, $43924444$ degree-two cocycle equations,
and every degree-two dimension and nonboundary generator independently.
All $123187$ zero-square certificates satisfy their literal coboundary
identity on $237022004$ normalized quadruples in total.  All $189$
remaining records have independently recomputed positive third homology.
The output is \path{code/outputs/b_fourteen_biset_primitives_verified.jsonl}.
The rational splittings used to skip projective-trivial opposite corners
are separately verified with exact fraction arithmetic in
\path{code/b_prepare_fourteen_biset_audit.py}.

The broader
order-nine resolution, stable-retract, rational-cup, and three/four-state
action calculations have already passed the separate checker
\path{code/b_verify_nine_small_acts.py}.  Its fresh all-element
action CSP independently matches all $18144$ residual action orbits
without using the original generator or stabilizer construction.
The global category lower bound is not advanced here: categories
with at least three objects require their own complete coverage.
\end{remark}
